\section{Simulation Bridge}\label{sec:min-validation}

The main application is semantic. The appendices keep controlled anchors where
the state is known in advance. These anchors show the patterns the audit is
designed to detect.

\paragraph{State, not scalar answers.}
HyperLogLog supplies the classical mergeable-sketch anchor. Cardinality itself
is underdetermined by $|A|$ and $|B|$ alone. The bounded register state
merges, and the noisy cardinality readout is applied after the state has been
composed. The appendix-side template carries over to the manifesto case. The
scalar economic score can remain nondecomposable across leaves; the richer
task state must be preserved so the root score stays readable.

\paragraph{Sufficiency under a channel.}
Sketch-Flip-Merge adds a privacy channel between local sketch construction and
the final merge \citep{HehirTingCormode2023SFM}. The lesson for C-TreePO is
not privacy itself but compatibility: a useful pipeline must preserve enough
state for the readout after every transformation it applies. If the local
randomizer, merge rule, and estimator are mismatched, the merged object no
longer lies in the family the readout expects. C-TreePO treats LLM
summarization as the semantic analogue and checks the compatibility condition
through C1/C2/C3 audits rather than a closed-form channel proof.

\paragraph{Ordered boundary state.}
The Markov/changepoint anchor is the clean version of a cross-span interaction
problem. A span state carries an internal count, a first register, and a last
register; merging two spans adds the two internal counts plus a boundary
indicator when the left last register differs from the right first register:
\[
  c=c_L+c_R+\One\{L.\mathrm{last}\ne R.\mathrm{first}\}.
\]
Smaller units can help when the merge state carries the right boundary
evidence. Lost boundary fields create the failure case. In the controlled
runs, local supervision can substitute for root labels because the local
labels measure the same oracle-preservation property. Appendix~\ref{app:v4-supervision-allocation}
expands the root/node supervision-budget evidence.

\paragraph{Learned realizations.}
Neural operators, transformers, embeddings, and prompted LLMs are different
ways to realize a local state map. The theorem stack is architecture-agnostic.
The implementation question is whether the realized family can approximate a
state function on the calls the tree actually makes, and whether the audit
shows that the approximation stayed inside an acceptable local-law residual.
Appendix~\ref{app:min-neural} keeps the formal neural-operator route;
Appendix~\ref{app:min-related} keeps broader representation and related-work
context.

The anchors support the same general pattern. If the state is rich enough for
the target and the local laws hold, the root is stable across merge schedules
and leaf sizes. If the state is too narrow, failures localize: leaves drop
evidence, merges lose boundary interactions, or re-summary drifts. The audit
serves the same diagnostic role in the Manifesto setting.
